% IEEE-styled paper. Uses the standard `article` class in two columns to remain
% portable; for a real submission, swap to \documentclass[conference]{IEEEtran}.
\documentclass[10pt,twocolumn]{article}
\usepackage[margin=0.75in]{geometry}
\usepackage{times}
\usepackage{booktabs}
\usepackage{amsmath}
\usepackage{caption}
\captionsetup{font=small,labelfont=bf}
\setlength{\parindent}{1em}
\title{\textbf{When mIoU Lies: Topology-Aware Evaluation of Thin Road-Marking Segmentation}}
\author{Anonymous Submission}
\date{}

\begin{document}
\maketitle

\begin{abstract}
Mean Intersection-over-Union (mIoU) is the de-facto metric for semantic
segmentation, including road-marking perception. We show two ways it fails for
\emph{thin} (2--8\,px) road markings. First, on \textbf{converged} models in a
width-controlled study (3 seeds, 95\% CIs), mIoU under-reports correctly and
continuously detected thin markings: at 2\,px, mIoU $=0.35\pm0.01$ while
centerline-Dice $=0.997\pm0.001$ and tolerance-IoU $=0.98$---a $\sim$0.65 gap with
non-overlapping confidence intervals. Second, under extreme multi-class imbalance
(real RLMD, 25 classes, 609:1), mIoU---in \emph{both} its all-class and foreground
forms---\emph{mis-ranks} losses: its top-ranked model is a frame-flooder (0.9\%
pixel accuracy) and it places the genuine best detector third. Frequency-weighted
IoU and pixel accuracy avoid rewarding flooding but are insensitive to foreground
quality (spread $\le 0.005$ vs.\ centerline-Dice's 0.31); only the topology metric
is both flood-robust and sensitive. We propose a thin-structure evaluation protocol
and an automated three-test check, and recommend reporting it alongside mIoU. A
studied connectivity loss yields no reliable multi-seed gain and is supplementary.
Results are on ground-view markings and a controlled synthetic benchmark; aerial
and full-scale validation are future work.
\end{abstract}

\section{Introduction}
Road-marking segmentation underpins lane perception, HD-map maintenance, and
infrastructure auditing. Markings are thin and sparse: in our real dataset they
occupy under 2\% of pixels, with a 609:1 imbalance between the most- and
least-common classes. The field reports mIoU almost universally. We ask whether
mIoU faithfully reflects thin-marking quality and selects the right model and
loss. It does not. Our contributions are: (1) a thin-structure evaluation protocol
(clDice, Boundary-IoU, tolerance-IoU, width-stratified analysis); (2) empirical
evidence that mIoU mis-ranks losses on real markings, crowning a model that floods
99\% of pixels, with the mechanism measured directly via pixel accuracy; (3) an
honest multi-seed study of a connectivity loss that yields no reliable gain; and
(4) a reproducible benchmark with a deployment path. This is an \emph{evaluation}
paper: we claim no new state of the art and report a loss that does not help.

\section{Related Work}
U-Net, U-Net++, DeepLabV3+, PSPNet, and SegFormer span the architecture space.
clDice with its differentiable soft-skeleton, and Boundary-IoU, were proposed for
tubular structures and boundary quality; we adopt them and add a width-stratified
view. Dice, Focal, and Tversky address imbalance; that overlap/accuracy metrics
reward over-prediction is known. Our contribution is the quantified rank-inversion
on real road markings, not the observation that imbalance is hard.

\section{Methodology}
\textbf{Protocol.} For binary any-marking masks we report clDice (harmonic mean of
skeleton precision/sensitivity), Boundary-IoU (IoU on a $d$-px boundary band),
tolerance-IoU (detection within $\tau$ px of the GT centerline), and a
width-stratified IoU/clDice where GT marking width is an independent variable.
\textbf{Loss.} We study $\mathcal{L}=w_{ce}\mathrm{CE}+w_r\mathrm{Dice}+
w_t(1-\text{soft-clDice})$. \textbf{Data.} RLMD (2{,}137 real images, 25 marking
classes) is the real benchmark; a controlled synthetic generator provides
ground-truth width control.

\section{Experimental Setup}
Single 4\,GB GPU; PyTorch; fp32. U-Net++ (ResNet-34) for the loss ablation.
ImageNet-pretrained encoders, AdamW, cosine. Compute-bounded: RLMD runs use a
500-image subset, 12 epochs, $384^2$. Statistics: per-image bootstrap 95\% CIs,
paired bootstrap tests, seeds 42/123/999. We disclose that the real 25-class
models are under-trained (best clDice 0.37).

\section{Results}
\subsection{mIoU under-reports thin markings (converged, multi-seed)}
On converged models (clDice $\approx0.997$ at the training width), three seeds,
evaluated across the controlled width series (Table~\ref{tab:width}; mean $\pm$
95\% CI): at 2\,px mIoU $=0.35\pm0.01$ while clDice $=0.997\pm0.001$ and
tolerance-IoU $=0.98$. The lines are correctly and continuously detected; the
$\sim$0.65 gap has non-overlapping, tight CIs and is \emph{not} an under-training
artifact. This is the paper's primary, hardest-to-attack result.

\begin{table}[t]\centering\small
\caption{Width study, converged models, mean $\pm$ 95\% CI over 3 seeds. mIoU
collapses on thin markings; topology/tolerance metrics do not.}
\label{tab:width}
\begin{tabular}{cccc}\toprule
width & mIoU & TolIoU & clDice \\\midrule
2\,px & $0.35\pm0.01$ & $0.98\pm0.01$ & $0.997\pm0.001$ \\
4\,px & $0.89\pm0.05$ & $1.00\pm0.00$ & $0.998\pm0.002$ \\
8\,px & $0.47\pm0.06$ & $0.97\pm0.01$ & $0.95\pm0.05$ \\\bottomrule
\end{tabular}
\end{table}

\subsection{mIoU mis-ranks losses (real RLMD)}
Table~\ref{tab:loss} gives the same-architecture loss ablation. We test every
common metric (Table~\ref{tab:tests}): does its top model flood the frame; does it
rank the genuine best detector first; is it foreground-sensitive? mIoU---in both
its all-class and foreground forms---ranks a frame-flooder first and the best
detector third; FWIoU and pixel accuracy avoid flooding but are insensitive
(spread $\le0.005$ vs.\ clDice's 0.31); only clDice passes all three. We scope this
as a multi-class extreme-imbalance phenomenon, complementing Sec.~V-A.

\begin{table}[t]\centering\small
\caption{RLMD loss ablation (U-Net++, same architecture). Tversky/Dice top mIoU
but flood the frame (PixelAcc $\approx 0.008$).}
\label{tab:loss}
\begin{tabular}{lccccc}\toprule
Loss & mIoU & PixAcc & clDice & clPrec & clRec \\\midrule
CE+Dice & 0.029 & 0.983 & \textbf{0.373} & 0.907 & 0.356 \\
CoRe & 0.005 & 0.981 & 0.275 & 0.896 & 0.211 \\
Focal & 0.018 & 0.983 & 0.085 & -- & 0.103 \\
CE & 0.010 & 0.982 & 0.063 & -- & 0.096 \\
Dice & 0.140 & 0.008 & 0.062 & -- & 1.000 \\
Tversky & \textbf{0.144} & \textbf{0.009} & 0.084 & 0.133 & 1.000 \\\bottomrule
\end{tabular}
\end{table}

\begin{table}[t]\centering\small
\caption{Three-test metric check on RLMD (extreme imbalance). Only clDice passes
all three.}
\label{tab:tests}
\begin{tabular}{lccc}\toprule
Metric & top floods? & best \#1? & spread \\\midrule
mIoU (fg) & yes & no (\#3) & 0.023 \\
mIoU (all) & yes & no (\#3) & 0.023 \\
FWIoU & no & yes & 0.005 \\
PixelAcc & no & yes & 0.003 \\
clDice & no & yes & \textbf{0.310} \\\bottomrule
\end{tabular}
\end{table}

\subsection{Connectivity loss: honest multi-seed result}
At seed 42, CE+Dice beats CoRe (clDice 0.37 vs.\ 0.28, $p<0.001$). Across three
seeds the advantage vanishes: CE+Dice $0.405\pm0.058$ vs.\ CoRe $0.406\pm0.141$
($\Delta=-0.001$, 95\% CI $[-0.18,0.18]$; CoRe wins at seed 999). The single-seed
result was an artifact caught only by multi-seed evaluation.

\subsection{Deployment}
ONNX Runtime is $1.47\times$ faster than eager PyTorch on CPU at bit-identical
output (max\,$|\Delta|=9\mathrm{e}{-8}$); a C++/LibTorch engine and TensorRT path
are released.

\section{Discussion and Limitations}
For thin road markings, report clDice/Boundary-IoU/tolerance-IoU and pixel
accuracy, not mIoU alone: mIoU can rank a useless frame-flooding model first.
\textbf{Limitations.} (1) RLMD is ground-view; we report no real aerial/UAV
marking results. (2) The multi-class real models are under-trained (clDice 0.37),
but the \emph{primary} result (Sec.~V-A) uses converged models (clDice $\approx
0.997$, 3 seeds, non-overlapping CIs) and a converged binary-marking model on real
RLMD plateaus at clDice 0.64 / mIoU 0.53$, so the central claim does not rest on
under-training. (3) The real mis-ranking is scoped to multi-class extreme
imbalance. (4) Three seeds; the cross-dataset probe is confounded; the studied loss
gives no reliable gain.

\section{Conclusion}
mIoU is an unreliable---sometimes rank-inverting---metric for thin road-marking
segmentation. We provide a topology-aware protocol that ranks models correctly, a
controlled demonstration of the mechanism, an honest multi-seed loss study, and a
reproducible benchmark. Converged real-data and aerial/UAV validation are the key
next steps.

\begin{thebibliography}{99}\small
\bibitem{unet} O. Ronneberger et al., ``U-Net,'' MICCAI, 2015.
\bibitem{unetpp} Z. Zhou et al., ``UNet++,'' DLMIA, 2018.
\bibitem{deeplab} L.-C. Chen et al., ``DeepLabV3+,'' ECCV, 2018.
\bibitem{psp} H. Zhao et al., ``PSPNet,'' CVPR, 2017.
\bibitem{segformer} E. Xie et al., ``SegFormer,'' NeurIPS, 2021.
\bibitem{cldice} S. Shit et al., ``clDice,'' CVPR, 2021.
\bibitem{biou} B. Cheng et al., ``Boundary IoU,'' CVPR, 2021.
\bibitem{dice} F. Milletari et al., ``V-Net,'' 3DV, 2016.
\bibitem{focal} T.-Y. Lin et al., ``Focal Loss,'' ICCV, 2017.
\bibitem{tversky} S. Salehi et al., ``Tversky loss,'' MLMI, 2017.
\bibitem{rlmd} ``RLMD: A Dataset for Road Marking Segmentation,'' IEEE.
\bibitem{loveda} J. Wang et al., ``LoveDA: A Remote Sensing Land-Cover Dataset for
Domain Adaptive Semantic Segmentation,'' NeurIPS D\&B, 2021.
\bibitem{uavid} Y. Lyu et al., ``UAVid: A Semantic Segmentation Dataset for UAV
Imagery,'' ISPRS J. Photogrammetry and Remote Sensing, 2020.
\bibitem{massroads} V. Mnih, ``Machine Learning for Aerial Image Labeling''
(Massachusetts Roads Dataset), PhD thesis, Univ.\ of Toronto, 2013.
\end{thebibliography}

\end{document}
